\section{Models and Objectives}\label{sec:models}

\subsection{Task 1: text tag classifier}
A DistilBERT encoder \cite{sanh2019distilbert}, a distilled version of BERT
\cite{devlin2019bert}, with a linear multi-label head predicts a 50-tag vocabulary from a
caption:
$t = \mathrm{BERT}_{\mathrm{CLS}}(X_{\mathrm{text}})$, $\hat{y}_k = \sigma(w_k^\top t + b_k)$,
trained with per-tag binary cross-entropy.

\subsection{Task 2: graph-based genre classifier and controls}
A 2-layer GraphSAGE encoder \cite{hamilton2017graphsage} with mean-pool readout classifies
GTZAN \cite{tzanetakis2002gtzan} genre, with each
non-final layer updating node representations as
\[
h_i^{(l+1)} = \sigma\big(W^{(l)} \cdot \mathrm{CONCAT}(h_i^{(l)}, \mathrm{MEAN}_{j\in\mathcal{N}(i)}h_j^{(l)})\big),
\]
(the implementation applies $\sigma$ and dropout between layers, not after the final one,
omitted above for brevity) before mean-pool readout $g = \tfrac{1}{|V|}\sum_i h_i^{(L)}$.
Three controls probe what the graph branch contributes, though none isolates topology
cleanly (Section~\ref{sec:protocol} discloses each control's specific confound): a CNN
baseline, following the standard convolutional auto-tagging approach for spectrogram-based
classification \cite{choi2016autotagging,won2020evaluation}, with no graph or shared
features; an edge-free set encoder replacing every GraphSAGE layer with a per-node MLP of
identical width/depth, keeping the same node features and readout but removing every edge
(changes the aggregation operator along with topology); and a non-learned exact 2-hop
neighbor aggregation, using GraphSAGE's topology and no trained message-passing function, but
a different training recipe for its head (Section~\ref{sec:discussion}). The cleaner
topology-only intervention this paper relies on is a same-architecture
real/shuffled/random-edge factorial from a companion audit \cite{tismir_audit}
(Section~\ref{sec:results-t2}), changing edges alone.

\subsection{Task 3: fusion strategies and diagnostics}
The same GNN and BERT component \emph{architectures} as Tasks 1 and 2 (not their trained
weights, separate freshly-initialized instances for Task 3) are combined four ways:
\textit{bert\_only} (text alone), \textit{gnn\_only} (graph alone), \textit{concat}
($z=\mathrm{CONCAT}(g,t)$), and \textit{cross\_attention}
($A=\mathrm{softmax}(QK^\top/\sqrt{d})$, $Q=gW_Q$, $K=H_{\mathrm{text}}W_K$,
$V=H_{\mathrm{text}}W_V$, $z=\mathrm{CONCAT}(g, AV)$; the implementation also adds a padding
mask to $QK^\top/\sqrt{d}$ before the softmax producing $A$, omitted above for brevity, so
padded tokens receive zero attention). Two diagnostics go beyond the four-way comparison.
\textbf{Branch zeroing}: inside a trained \textit{concat} model, the graph embedding $g$ is
replaced with zeros at evaluation time, no retraining, isolating what the graph branch
contributes given text is already present. \textbf{Complementarity cross-tabulation}:
per-(example, tag) predictions from \textit{bert\_only} and \textit{gnn\_only} are compared
directly, testing whether the graph model is right on cases text gets wrong, rather than only
comparing aggregate scores.

\subsection{Task 4: contrastive dual encoder and controls}
A CLIP-style \cite{radford2021clip} dual encoder pairs the same GNN and BERT components under
a symmetric InfoNCE objective \cite{oord2018infonce}. The same two-tower, contrastive design
underlies music-specific audio-text systems -- MuLan \cite{huang2022mulan} and Manco et
al.'s contrastive audio-language model \cite{manco2022contrastive} -- and general-purpose
audio-language systems, such as CLAP \cite{elizalde2023clap}; the same objective is used for
text-to-music retrieval by Doh et al.~\cite{doh2023universal}. This task follows that design
at a much smaller scale, with both towers built from this paper's own GNN and BERT components
rather than larger pretrained encoders -- averaging the audio-to-text direction with
text-to-audio (swapping which side sums in the denominator) gives the loss $\mathcal{L} =
\tfrac{1}{2}(\mathcal{L}_{\mathrm{a2t}} + \mathcal{L}_{\mathrm{t2a}})$ its "symmetric" name;
one direction shown for brevity:
\[
\mathcal{L}_{\mathrm{a2t}} = -\frac{1}{N}\sum_i \log\frac{\exp(\mathrm{sim}(g_i,t_i)/\tau)}{\sum_j \exp(\mathrm{sim}(g_i,t_j)/\tau)}.
\]
The same graph-vs-edge-free comparison from Task 2 is repeated here: an otherwise identical
dual encoder with the GNN branch replaced by a width/depth-matched, lower-parameter, edge-free MLP encoder,
trained and evaluated identically.
